\section{Proof of Theorem H}\label{sec:hypothesis}

This section is not formalized. We keep the notation of \cref{sec:operator}: $E=\Q(\sqrt5)$, $\OO_E=\Z[\theta]$, the
icosian ring $I$, the group $\Gamma^+=I[1/3]^1/\{\pm1\}$ acting on the tree $\Tree$ of $\PGL_2(\Q_9)$, and the prime
$\lambda_0\mid11$. We use that $\theta$ has norm $-1$, that $E$ has narrow class number one, and that every totally
positive unit of $\OO_E$ is a square.

Theorem~H (\cref{thm:H}) is proved in four steps. (H1)~The group $\Gamma^+$ is an amalgam of two copies of~$A_5$
(\cref{lem:A5,lem:amalgam}), its congruence subgroups form a chain (\cref{lem:chain}), and the kernels of $A$ along that
chain are Tanner codes of quotients of~$\Tree$ (\cref{lem:tanner,lem:covering}). (H2)~Riemann--Roch on a Hilbert modular
surface in characteristic~$3$ gives many sections of a line bundle on its Goren--Oort divisor (\cref{lem:RR}). (H3)~That
divisor is a quotient of a union of projective lines indexed by~$\Tree$, on which the line bundle has degree~$4$
(\cref{lem:uniformization}), so its sections form a Tanner code (\cref{lem:sections}). (H4)~This Tanner code is the one of
the twisted datum (\cref{lem:rigidity,lem:twist}). Steps (H2) and~(H3) rest on published theorems, listed where they are
used. The new arguments are \cref{lem:rigidity}, the degree computation in \cref{lem:uniformization}, \cref{lem:twist},
and the assembly in \cref{sec:proofH}.

\subsection{The group, the chain and the covering}\label{sec:groupchain}

Let $\sigma$ be the Frobenius of $\F_9$, and also its lift to~$\Q_9$.

\begin{lemma}[$A_5$ on the projective line]\label{lem:A5}
\notformalized
Let $A_5$ act faithfully on $\Pone$ over $\bar\F_3$.
\begin{enumerate}[label=(\alph*)]
\item The action comes from an irreducible faithful representation of degree~$2$ of the Schur cover $2I\cong\SL_2(\F_5)$
on which $-1$ acts by~$-1$. Up to isomorphism there are exactly two such representations; they are defined over $\F_9$,
are not isomorphic, and are exchanged by~$\sigma$.
\item The centralizer of $A_5$ in $\PGL_2(\bar\F_3)$ is trivial.
\item The orbits of $A_5$ on $\Pone(\bar\F_3)$ have sizes $10$, $12$ and $60$ only. The orbit of size $10$ consists of the
fixed points of the ten subgroups $S_3$ of~$A_5$; each of them fixes exactly one point, and $S_3$ is self-normalizing
in~$A_5$.
\end{enumerate}
\end{lemma}
\begin{informalproof}
(a) A faithful action of $A_5$ on $\Pone$ is a faithful projective representation of degree~$2$. The group $A_5$ has no
faithful representation of degree~$2$ in characteristic~$3$ (its $3$-modular irreducible representations have degrees
$1,3,3,4$~\cite{JLPW1995}), so the action lifts to a faithful representation of $2I$ on which $-1$ acts by~$-1$. It is
irreducible, since a reducible one has image in a Borel subgroup, which is solvable. The faithful $3$-modular irreducible
representations of $\SL_2(\F_5)$ have degrees $2,2,6$~\cite{JLPW1995}, so there are exactly two of degree~$2$. Their
traces on an element of order~$5$ are the two roots of $x^2+x-1$ in~$\F_9$, so they are defined over $\F_9$, not
isomorphic, and exchanged by~$\sigma$.

(b) If $g\in\PGL_2(\bar\F_3)$ commutes with the image of $A_5$, and $\rho$ is a representation as in~(a), then
$g\rho(h)g^{-1}=c(h)\rho(h)$ for a character $c$ of $2I$, which is trivial because $2I$ is perfect; so $g$ commutes with
the absolutely irreducible $\rho(2I)$ and is a scalar.

(c) The stabilizer of a point lies in a Borel subgroup, whose finite subgroups are extensions of a cyclic group of order
prime to~$3$ by a $3$-group. Among the subgroups of $A_5$ only $1$, $C_2$, $C_3$, $C_5$ and $S_3$ have this form. An
element of order~$3$ is unipotent and has one fixed point, which is fixed by its normalizer~$S_3$; so each $S_3$ fixes
exactly that point, and two different subgroups $S_3$ fix different points, because together they generate~$A_5$, which
fixes no point. The ten subgroups $S_3$ are conjugate, so their fixed points form one orbit of size~$10$. Each of these
points is fixed by the three involutions of its~$S_3$, which accounts for all $15\cdot2$ fixed points of the fifteen
involutions; so no point has stabilizer~$C_2$. An element of order~$5$ is semisimple with two fixed points, whose
stabilizer is $C_5$, since the dihedral group of order~$10$ is not of the form above. This leaves orbits of sizes $10$,
$12$ and~$60$. Finally $S_3$ is maximal and not normal in~$A_5$, hence self-normalizing.
\end{informalproof}

\begin{lemma}[The amalgam]\label{lem:amalgam}
\notformalized
$\Gamma^+$ acts on $\Tree$ without inversions, transitively on the even vertices, on the odd vertices and on the edges.
The stabilizers of $\Lambda_0$, $\Lambda_1$ and $e_0$ are $A_5=2I/\{\pm1\}$, $\alpha A_5\alpha^{-1}$ and a subgroup $S_3$,
whose preimage $\Gamma_e\subset2I$ is binary dihedral of order~$12$. Hence $\Gamma^+=A_5*_{S_3}\alpha A_5\alpha^{-1}$, and
$\Gamma^+$ is perfect.
\end{lemma}
\begin{informalproof}
Elements of reduced norm $1$ preserve the parity of vertices. A vertex of $\Tree$ is the class of a $\Z_9$-lattice in
$\Q_9^2$; by the local--global correspondence for lattices~\cite[Section~III.5.A]{Vigneras1980} it is the class of the
completion at~$3$ of a right $I$-ideal $\mathfrak a$ with $\mathfrak a_{\mathfrak q}=I_{\mathfrak q}$ for all primes
$\mathfrak q\nmid3$, and such an $\mathfrak a$ is principal, $\mathfrak a=\beta I$ with $\beta\in I[1/3]^\times$, because
$I$ has class number one (\cite[p.~155]{Vigneras1980}, the case $m=5$, $D=N=1$ of the table there). The reduced norm of
$\beta$ is a totally positive unit of $\OO_E[1/3]$, so it is $3^a\theta^{2b}$; replacing $\beta$ by $\beta\theta^{-b}$ we may
take $b=0$. If $a$ is even, $\beta/3^{a/2}\in\Gamma^+$ moves $\Lambda_0$ to the vertex of~$\mathfrak a$; if $a$ is odd,
$(\beta/3^{(a-1)/2})\alpha^{-1}$ moves $\Lambda_1$ to it. So $\Gamma^+$ is transitive on the vertices of each parity. The
stabilizer of $\Lambda_0$ is the image of $I^1=2I$. Reduction modulo~$3$ is injective on the finite group $2I$, because
$1+3M_2(\Z_9)$ is torsion-free, so $A_5$ acts faithfully on the $10$ points of $\Pone(\Lambda_0/3\Lambda_0)=\Pone(\F_9)$,
which correspond to the edges at~$\Lambda_0$. By \cref{lem:A5}(c) they form one orbit, with stabilizer $S_3$, whose
preimage in $2I\cong\SL_2(\F_5)$ is binary dihedral of order~$12$. Hence $\Gamma^+$ is transitive on edges, and Bass--Serre
theory~\cite[Section~I.4]{Serre1980} gives the amalgam. As a check, the Euler characteristic $22-110=-88$ of $X_0$ equals
$660\,(1/60+1/60-1/6)$. An amalgam of perfect groups is perfect.
\end{informalproof}

\begin{lemma}[The chain]\label{lem:chain}
\notformalized
For $m\ge1$ the map $\Gamma^+\to\PSL_2(\OO_E/\lambda_0^m)=\PSL_2(\Z/11^m)$ is onto, and $\Gamma^+(\lambda_0^m)$ is a
torsion-free normal subgroup of $\Gamma^+$ that acts freely on~$\Tree$. The groups $\Gamma^+(\lambda_0^{k+1})$,
$k\ge0$, are strictly decreasing, $\abs{\Gamma^+(\lambda_0):\Gamma^+(\lambda_0^{k+1})}=11^{3k}$, and their intersection is
trivial. Moreover $\Gamma^+(\lambda_0^m)\cong\{\gamma\in I[1/3]^1:\gamma\equiv1\pmod{\lambda_0^m}\}$.
\end{lemma}
\begin{informalproof}
The algebraic group $B^1$ of norm-one elements is simply connected and absolutely almost simple, and $B^1(E_3)=\SL_2(\Q_9)$
is not compact. Strong approximation~\cite[Section~7.4]{PlatonovRapinchuk1994} makes $I[1/3]^1$ dense in
$\prod_{\mathfrak q\nmid3}(I\otimes\OO_{E,\mathfrak q})^1$, so $I[1/3]^1\to\SL_2(\OO_E/\lambda_0^m)$ is onto. The kernel
$\Gamma^+(\lambda_0^m)$ is normal. A nontrivial element of finite order fixes a vertex of~$\Tree$~\cite[Section~I.4]{Serre1980},
so it lies in a $\Gamma^+$-conjugate of $A_5$ or of $\alpha A_5\alpha^{-1}$; these map injectively to
$\PSL_2(\F_{11})$, because reduction modulo~$\lambda_0$ is injective on $2I$ ($1+\lambda_0M_2(\OO_{E,\lambda_0})$ is
torsion-free). So $\Gamma^+(\lambda_0^m)$ is torsion-free, and since $\Gamma^+$ acts without inversions, it acts freely.
Since $-1\not\equiv1\pmod{\lambda_0}$, each class $\pm\gamma$ in $\Gamma^+(\lambda_0^m)$ has exactly one representative
$\gamma\equiv1\pmod{\lambda_0^m}$. The index is the order of $\ker(\PSL_2(\Z/11^{k+1})\to\PSL_2(\F_{11}))$, which is
$11^{3k}$. If $g\in I[1/3]^1$ satisfies $g\equiv\pm1\pmod{\lambda_0^{k+1}}$ for every~$k$, the sign is constant because
$1\not\equiv-1\pmod{\lambda_0}$, so $g=\pm1$ in $I\otimes\OO_{E,\lambda_0}$, and $g=\pm1$ because $B\to B\otimes E_{\lambda_0}$ is
injective.
\end{informalproof}

\begin{lemma}[Covering theory]\label{lem:covering}
\notformalized
There is an isomorphism $\psi\colon F=\pi_1(X_0,\tau)\to\Gamma^+(\lambda_0)$ such that for every normal subgroup $N$ of $F$
of finite index the cover $X_N$ of \cref{lem:tanner}, with its lifted datum, is isomorphic to $\psi(N)\backslash\Tree$ with
the datum descended from $D_{\mathrm{tw}}$. Consequently
$\dim_{\F_3}\ker\varphi_{F/N}(A)=2\dim_{\F_9}T_{\mathrm{tw}}(\psi(N)\backslash\Tree)$.
\end{lemma}
\begin{informalproof}
Since $\Gamma^+(\lambda_0)$ acts freely and without inversions, $\Tree\to X_0$ is a universal covering, and lifting the
spanning tree $\tau$ to~$\Tree$ gives an isomorphism $\psi$ that sends the generator of an edge $e$ to the deck
transformation of the fundamental cycle of~$e$. The identification must match the data, not only the graphs, since $X_0$
has many automorphisms. We use the identification of \cref{lem:basegraph}, through which the data of \cref{def:operator}
are defined: labels are transported by coset representatives, so the datum of $X_0$ is the descent of $D_{\mathrm{tw}}$.
Then $X_N$ and $\psi(N)\backslash\Tree$ carry the descent of the same datum. The last formula is \cref{lem:tanner}.
\end{informalproof}

Let $N_k=\psi^{-1}(\Gamma^+(\lambda_0^{k+1}))$ for $k\ge0$. By \cref{lem:chain}, $(N_k)$ is a chain of $F$ with
$\abs{F:N_k}=11^{3k}$, and with $X_k=\Gamma^+(\lambda_0^{k+1})\backslash\Tree$, \cref{lem:covering} gives
\begin{equation}\label{eq:H1}
  \dim_{\F_3}\ker\varphi_{F/N_k}(A)=2\dim_{\F_9}T_{\mathrm{tw}}(X_k).
\end{equation}

\subsection{Riemann--Roch on the Hilbert modular surface}\label{sec:RR}

Let $\mathfrak n=\lambda_0^{k+1}$, an ideal prime to~$3$ with a totally positive generator, and let
$\mathcal K(\mathfrak n)=\ker\bigl(\GL_2(\hat\OO_E)\to\GL_2(\OO_E/\mathfrak n)\bigr)$ be the principal congruence subgroup
of level~$\mathfrak n$. It is neat: a root of unity congruent to~$1$ modulo a prime above~$11$ has $11$-power order, so it
is $1$ or has degree at least $10$ over~$\Q$, while the eigenvalues of elements of $\GL_2(E)$ have degree at most~$4$. Choose a
smooth projective toroidal compactification, over $W(\bar\F_3)$, of the closure of one geometric connected component of the
Hilbert modular surface of $E$ of level $\mathcal K(\mathfrak n)$, compatibly as $k$ varies, and let $S_k$ be its special
fibre. The complex fibre of that component is
$\Gamma(\mathfrak n)\backslash\mathfrak H^2$ with $\Gamma(\mathfrak n)\subset\SL_2(\OO_E)$, since $E$ has narrow class number
one and every totally positive unit is a square. We use the following published facts. The sources state them for full level $n\ge3$, for
$\Gamma_1(\mathfrak c,\mathfrak n)$-structures or for general neat levels; the facts we use are local on the moduli space or
descend along the finite \'etale maps between levels prime to~$3$, so they hold at level~$\mathcal K(\mathfrak n)$.
\begin{itemize}
\item[(G1)] $S_k$ is smooth, projective and connected. Its boundary $D_k$ is a disjoint union of $c_k\ge1$ cycles of smooth
rational curves, one for each cusp~\cite[Theorem~5.1, Proposition~5.2 and Corollary~5.3]{Rapoport1978}, \cite{Chai1990},
\cite[Section~7]{Dimitrov2004}, \cite[Theorems~6.4.1.1 and~7.3.3.4]{Lan2013}.
\item[(G2)] Since $3$ is unramified in~$E$, $\OO_E\otimes W=W\times W$ and the Hodge bundle splits as
$\omega_1\oplus\omega_2$. The Kodaira--Spencer map with logarithmic poles gives
$\Omega^1_{S_k}(\log D_k)\cong\omega_1^2\oplus\omega_2^2$~\cite[Theorem~6.4.1.1(4)]{Lan2013}: Lan's Kodaira--Spencer sheaf is the
symmetric square of $\omega$ over $\OO_E\otimes\OO$, because $\omega$ is locally free of rank one over $\OO_E\otimes\OO$ and a
polarization of degree prime to~$3$ identifies the Hodge bundle of the dual with~$\omega$. Hence
\[
  K_{S_k}+D_k\cong\omega_1^2\otimes\omega_2^2=:L,
\]
the bundle of parallel weight~$2$.
\item[(G3)] The partial Hasse invariants $h_1\in H^0(\omega_1^{-1}\omega_2^3)$ and $h_2\in H^0(\omega_1^3\omega_2^{-1})$ have
reduced divisors $Z_1$, $Z_2$, the Goren--Oort divisors, which are regular curves meeting in finitely many
points~\cite[Section~1 and Theorem~2.1(1)]{Goren2001}; see also~\cite{GorenOort2000,TianXiao2016}. They extend to the
compactification by the Koecher principle~\cite[Theorem~8.3]{Dimitrov2004}, and they do not vanish on~$D_k$, because their
$q$-expansion at every cusp is~$1$~\cite[Theorem~2.1(2)]{Goren2001}. So $h_1h_2\in H^0(L)$ has divisor $Y_k=Z_1+Z_2$ and
$L(-Y_k)\cong\OO$. Since $3$ is inert, $Y_k$ is the reduced supersingular locus~\cite[Proposition~4.1]{BachmatGoren1999}.
\item[(G4)] $\chi(\OO_{S_k})=V_k:=e_{\mathrm{top}}/4=11^{3k+1}$, where $e_{\mathrm{top}}$ is the Euler number of
$\Gamma(\mathfrak n)\backslash\mathfrak H^2$. Since $\chi(\OO)$ is constant in flat families and a birational invariant, it
may be computed in characteristic~$0$ with any resolution of the cusps. Every cusp of $\Gamma(\mathfrak n)$ is conjugate
to~$\infty$ and has type $(\mathfrak n,\mathcal U_{\mathfrak n})$, which is isomorphic to $(\OO_E,\mathcal U_{\mathfrak n})$
through a totally positive generator of~$\mathfrak n$, where
$\mathcal U_{\mathfrak n}=\{\eps^2:\eps\in\OO_E^\times,\ \eps\equiv1\pmod{\mathfrak n}\}$. The module $\OO_E=\Z+\Z\theta^2$ has
reduced number $\theta^2=3-1/\theta^2$, with purely periodic minus continued fraction $[[3]]$, so every cusp is resolved by a
cycle of $[\OO_E^{\times,+}:\mathcal U_{\mathfrak n}]$ curves of self-intersection~$-3$~\cite[Sections~2.5 and~2.6]{Hirzebruch1973},
\cite{vanderGeer1988}:
$5$ curves for $k=0$, since $\theta\equiv4$ has order~$5$ modulo~$\lambda_0$, and $55$ for $k=1$, since $\theta\equiv37$ has
order~$55$ modulo~$\lambda_0^2$. Noether's formula, with $K^2=2e_{\mathrm{top}}+\sum_j(2-b_j)$ and
$c_2=e_{\mathrm{top}}+\#\{\text{curves}\}$ for curves of self-intersection $-b_j=-3$, gives
$\chi=(K^2+c_2)/12=e_{\mathrm{top}}/4+\sum_j(3-b_j)/12=e_{\mathrm{top}}/4$. Finally $e_{\mathrm{top}}=\abs{\PSL_2(\F_{11})}\cdot
2\zeta_E(-1)=660/15=44$ at level $\lambda_0$, since $2\zeta_E(-1)=1/15$~\cite[Sections~1.6 and~1.7]{Hirzebruch1973}, and $e_{\mathrm{top}}$ is multiplied by $11^{3k}$ at
level $\lambda_0^{k+1}$.
\end{itemize}

\begin{lemma}[Sections on the Goren--Oort divisor]\label{lem:RR}
\notformalized
$h^0(Y_k,L|_{Y_k})\ge V_k-1$.
\end{lemma}
\begin{informalproof}
By Riemann--Roch, $\chi(L)=\chi(\OO)+\frac12L\cdot(L-K)=\chi(\OO)+\frac12(K+D)\cdot D$. Each boundary cycle $D_j$ has
arithmetic genus~$1$ and the $D_j$ are disjoint, so adjunction gives $(K+D)\cdot D=\sum_j(K+D_j)\cdot D_j=0$. Hence
$\chi(L)=\chi(\OO)=V_k$ by (G4). By Serre duality $h^2(L)=h^0(K-L)=h^0(-D)=0$, because $D\ne0$ and $S_k$ is integral.
So $h^0(L)\ge\chi(L)=V_k$. The exact sequence $0\to L(-Y_k)\cong\OO\to L\to L|_{Y_k}\to0$ and $h^0(\OO)=1$ give
$h^0(Y_k,L|_{Y_k})\ge h^0(L)-1\ge V_k-1$.
\end{informalproof}

If moreover $h^1(\OO_{S_k})=0$, then $h^1(L)=h^1(-D)=c_k-1$ and $h^0(L|_{Y_k})=V_k+c_k-2$ exactly; the proof uses only the
lower bound.

\subsection{Uniformization of the supersingular locus}\label{sec:uniformization}

We use the description of the supersingular locus of Hilbert--Blumenthal surfaces at an inert prime by Stamm, as
presented by Kudla and Rapoport~\cite[Section~4]{KudlaRapoport1999}, \cite{Stamm1997}, the uniformization theorem of
Rapoport and Zink~\cite[Theorem~6.30]{RapoportZink1996}, and the intersection behaviour of the Goren--Oort divisors found
by Bachmat and Goren~\cite[Theorem~6.1]{BachmatGoren1999}. Kudla and Rapoport work with abelian varieties of
dimension~$8$ with an action of $C(V)\otimes k$, for the Clifford algebra $C(V)$ of a quadratic space attached to a
quaternion algebra $B_0$ over~$\Q$; for $B_0=M_2(\Q)$ and $k=E$ their moduli space is the Hilbert--Blumenthal surface,
and Morita equivalence identifies their Dieudonn\'e lattices with ours~\cite[Introduction and Section~3]{KudlaRapoport1999}.

\paragraph{The local picture.} Let $\mathcal M$ be the Rapoport--Zink space of supersingular polarized $3$-divisible
groups with $\OO_E\otimes\Z_3$-action, with framing object~$\mathbb X$: a point is such a group $X$ with a framing
quasi-isogeny $\varrho$ between $X$ and~$\mathbb X$. The group $J$ of $\OO_E$-linear self-quasi-isogenies of $\mathbb X$
that multiply the polarization by a factor in $\Q_3^\times$ acts on $\mathcal M$ by changing~$\varrho$, so the Hodge bundles
$\omega_1$, $\omega_2$ of the universal $3$-divisible group, and $L=\omega_1^2\otimes\omega_2^2$, carry a canonical
$J$-action. Following~\cite[Section~4]{KudlaRapoport1999}, the isocrystal of $\mathbb X$ yields two-dimensional
$\Q_9$-spaces $U_0$ and $U_1$ with alternating forms $\langle\,,\rangle$, on which $J$ acts, and Frobenius gives a
$\sigma$-semilinear bijection $\Phi\colon U_1\to U_0$ that commutes with~$J$; Frobenius maps $U_0$ to $U_1$ by
$3\Phi^{-1}$, and $\langle\Phi x,\Phi y\rangle=3\sigma\langle x,y\rangle$. For $n\in\Z$ let $\mathcal M^{(n)}_{\mathrm{red}}$ be the
open and closed part of the reduced locus whose Dieudonn\'e lattices $L_0\oplus L_1$ satisfy $\langle L_i,L_i\rangle=3^nW$;
its stabilizer $J^0$ in $J$ consists of the elements whose factor is a $3$-adic unit.
By~\cite[Lemma~4.3 and Proposition~4.4]{KudlaRapoport1999}, $\mathcal M^{(n)}_{\mathrm{red}}$ is a union of projective lines
over~$\bar\F_3$: a line $C_\Lambda=\Pone(\Lambda/3\Lambda)$ of
the first kind for each $\Z_9$-lattice $\Lambda\subset U_0$ with $\langle\Lambda,\Lambda\rangle=3^n\Z_9$, and a line
$C_{\Lambda'}=\Pone(\Lambda'/3\Lambda')$ of the second kind for each $\Z_9$-lattice $\Lambda'\subset U_1$ with
$\langle\Lambda',\Lambda'\rangle=3^n\Z_9$. Lines of the same kind are disjoint. The lines $C_\Lambda$ and $C_{\Lambda'}$ meet
if and only if $3\Lambda\subset\Phi\Lambda'\subset\Lambda$, and then in one point, which is $\Phi\Lambda'/3\Lambda$ on
$C_\Lambda$ and $\Phi^{-1}(3\Lambda)/3\Lambda'$ on $C_{\Lambda'}$~\cite[Lemma~4.3(ii)]{KudlaRapoport1999}. Each line meets
exactly one line of the other kind at each of its ten $\F_9$-rational points, the superspecial points, and no line
elsewhere~\cite[Lemma~4.3(iii) and the text after~(4.9)]{KudlaRapoport1999}. As schemes over~$\F_9$, the lines are the
projective lines $\Pone(\Lambda/3\Lambda)$ and $\Pone(\Lambda'/3\Lambda')$, with quasi-isogenies acting through the lattices:
the morphism from their union to the supersingular locus is its normalization, equivariantly~\cite[(4.10)]{KudlaRapoport1999},
after Stamm~\cite{Stamm1997}. The crossings are ordinary double points~\cite[Theorem~6.1]{BachmatGoren1999}. These two
statements concern the global moduli space, and they transfer to $\mathcal M$ through the uniformization theorem used below;
the proof of the second is local deformation theory. Scaling a lattice by $3$
multiplies $\langle\,,\rangle$ by~$9$, and $\langle\Phi\Lambda',\Phi\Lambda'\rangle=3^{n+1}\Z_9$; so $C_\Lambda\mapsto[\Lambda]$
and $C_{\Lambda'}\mapsto[\Phi\Lambda']$ map the lines of the two kinds bijectively to the homothety classes of lattices in
$U_0$ of the two parities, and they identify the dual graph of $\mathcal M^{(n)}_{\mathrm{red}}$ with the tree
$\Tree_{U_0}$ of~$U_0$, equivariantly for~$J^0$.

\begin{lemma}[Uniformization]\label{lem:uniformization}
\notformalized
There exist an integer~$n$, an action of $\Gamma^+$ on $\mathcal M^{(n)}_{\mathrm{red}}$ through~$J^0$, and an isomorphism
$U_0\cong\Q_9^2$ that carries $\Tree_{U_0}$ to~$\Tree$, with the following properties.
\begin{enumerate}[label=(\alph*)]
\item The resulting identification of the dual graph of $\mathcal M^{(n)}_{\mathrm{red}}$ with $\Tree$ is equivariant for
the action of $\Gamma^+$ on $\Tree$ of \cref{sec:operator}.
\item For every $k\ge0$, $Y_k\cong\Gamma^+(\lambda_0^{k+1})\backslash\mathcal M^{(n)}_{\mathrm{red}}$, and under this
isomorphism $L|_{Y_k}$ is the descent of $L$ with its $J$-action. In particular the dual graph of $Y_k$ is~$X_k$.
\item $L$ has degree~$4$ on every line of $\mathcal M^{(n)}_{\mathrm{red}}$.
\end{enumerate}
\end{lemma}
\begin{informalproof}
\emph{The global picture.} Let $G'=\{g\in B^\times:\Nrd g\in\Q^\times\}$, write $\mathcal K(\mathfrak n)=\mathcal
K_3\mathcal K^{(3)}$ with $\mathcal K^{(3)}$ the part away from~$3$, and let $\mathbb A_f^{(3)}$ be the finite adeles away
from~$3$. Fix a supersingular base point $x_0$ of $S_0$ with integral level structure, which exists by \cref{lem:RR},
with compatible lifts to the~$S_k$, and frame $\mathcal M$ by its $3$-divisible group. The supersingular points form one
isogeny class, the class~$I_0$ of~\cite[Section~3]{KudlaRapoport1999}, which is basic. The uniformization
theorem~\cite[Theorem~6.30]{RapoportZink1996} identifies the formal completion of the Shimura variety of level
$\mathcal K(\mathfrak n)$ along an open and closed part $Z'$ of the supersingular locus $Z$ with
$G'(\Q)\backslash(\mathcal M\times G'(\mathbb A_f^{(3)})/\mathcal K^{(3)})$, a finite disjoint union of quotients of~$\mathcal M$.
Here $G'(\Q)$ is the group of $\OO_E$-linear self-quasi-isogenies of~$x_0$ that multiply the polarization by a factor
in~$\Q^\times$; it acts on $\mathcal M$ through~$J$, and on $G'(\mathbb A_f^{(3)})$ through the Tate modules of $x_0$ and its
level structure. Moreover $Z'=Z$, because the group $G^*=\{g\in\GL_2(E):\det g\in\Q^\times\}$ of the moduli problem
satisfies the Hasse principle: its derived group $\operatorname{Res}_{E/\Q}\SL_2$ is simply connected, so it satisfies the
Hasse principle~\cite[Chapter~6]{PlatonovRapinchuk1994}; $G^*/G^{*,\mathrm{der}}\cong\mathbb G_m$ has trivial first Galois
cohomology; and $\det\colon G^*(\Q_v)\to\Q_v^\times$ is onto at every place~$v$. The endomorphisms of~$x_0$, with $3$
inverted, form a maximal $\OO_E[1/3]$-order of the definite quaternion algebra~$B$, by Tate's theorem~\cite{Tate1966} at
the primes $\ell\ne3$. Since $B$ has class number one, it has type number one~\cite[Corollaire~I.4.11 and p.~155]{Vigneras1980},
and we may choose the identification of $B$ with the endomorphism algebra of $x_0$ so that this order is $I[1/3]$. By the norm theorem of Hasse and Schilling, which gives $\Q_{>0}\subseteq\Nrd B^\times$~\cite[Theorem~III.4.1]{Vigneras1980},
and strong approximation for $B^1$ with respect to~$\{3\}$, every supersingular point of~$S_k$ has a representative in the slice $\mathcal M^{(n)}_{\mathrm{red}}\times\{\mathcal K^{(3)}\}$ that contains~$x_0$. The
stabilizer of that slice is $\{\gamma\in G'(\Q)\cap\mathcal K^{(3)}:v_3(\Nrd\gamma)=0\}$. Its reduced norms are positive
rationals that are units at every prime, so they equal~$1$, and the stabilizer is
$\{\gamma\in I[1/3]^1:\gamma\equiv1\pmod{\mathfrak n}\}\cong\Gamma^+(\lambda_0^{k+1})$ (\cref{lem:chain}). The element
$-1\in I[1/3]^1$ acts trivially on $\mathcal M_{\mathrm{red}}$, since $(X,\varrho)$ is isomorphic to $(X,-\varrho)$ through
$-1\in\operatorname{Aut}X$, and on~$L$, by $(-1)^4=1$; so $I[1/3]^1$ acts through $\Gamma^+$, and through $J^0$ since its
reduced norms are~$1$. The uniformization is compatible with the universal $3$-divisible groups, hence with $\omega_1$,
$\omega_2$ and~$L$. This proves~(b), given~(a).

\emph{The identification with $\Tree$.} The algebra $B$ acts on the isocrystal of $x_0$ contravariantly. Composed with
the main involution $b\mapsto\bar b$, this action gives a homomorphism $B\otimes_EE_3\to\operatorname{End}_{\Q_9}(U_0)$, an
isomorphism of central simple algebras, which on elements of reduced norm one is their action on lattices. Here $E_3$ is
identified with $\Q_9$ through the embedding of $\OO_E\otimes\Z_3$ into $W$ that defines~$L_0$. By the Skolem--Noether
theorem we may choose $U_0\cong\Q_9^2$ so that this homomorphism becomes the isomorphism $B\otimes_EE_3\cong M_2(\Q_9)$ of
\cref{sec:operator}, composed with $\sigma$ if that embedding reduces to $\theta\mapsto2-i$ (a case treated in
\cref{lem:twist}); then the identification of the dual graph with $\Tree_{U_0}\cong\Tree$ is $\Gamma^+$-equivariant, which
is~(a).

\emph{The degree.} Since $\Gamma^+(\lambda_0^{k+1})$ acts freely on $\Tree$ (\cref{lem:chain}), every line maps
isomorphically onto a component of~$Y_k$, compatibly with~$L$. Let $C$ be a component of $Z_1$, and put
$a=\deg\omega_1|_C$ and $b=\deg\omega_2|_C$. The section $h_2$ vanishes on $C$ exactly at the ten superspecial points, simply,
because the crossings are ordinary double points of the reduced curve $Z_1+Z_2$~\cite[Theorem~6.1]{BachmatGoren1999}; so
$3a-b=10$. Lines of the same kind are disjoint, so $C^2=Z_1\cdot C=-a+3b$, and adjunction with $K=L-D$ and $D\cdot C=0$
gives $-2=(K+C)\cdot C=a+5b$. Hence $(a,b)=(3,-1)$ and $\deg L|_C=2a+2b=4$; the components of $Z_2$ are symmetric. This
is~(c). As a check, $Z_1\sim3\omega_2-\omega_1$, $Z_2\sim3\omega_1-\omega_2$ and $\omega_i^2=0$ give
$Z_1\cdot Z_2=10\,\omega_1\omega_2=10\,e_{\mathrm{top}}/4$, while each of the $n_1$ lines of one kind meets ten lines of the
other, so $Z_1\cdot Z_2=10n_1$; thus $n_1=e_{\mathrm{top}}/4=11^{3k+1}=\abs{\Gamma^+:\Gamma^+(\lambda_0^{k+1})}/60$, as $X_k$
requires.
\end{informalproof}

\paragraph{Geometric data.} A \emph{geometric datum} on $\Tree$ is a $\Gamma^+$-equivariant family consisting of a smooth
rational curve $C_v$ over $\bar\F_3$ for each vertex~$v$, on which the stabilizer of $v$ acts faithfully; a line bundle
$L_v$ of degree~$4$ on~$C_v$, with an action of that stabilizer; for each edge $e$ at~$v$ a point $P_{v,e}\in C_v$, the
\emph{special points}; and for each edge $e=\{v,w\}$ an isomorphism between the fibres of $L_v$ at $P_{v,e}$ and of $L_w$
at~$P_{w,e}$. Its \emph{type} is the pair of the isomorphism classes of the actions of the stabilizers of $\Lambda_0$ and
$\Lambda_1$ on $C_{\Lambda_0}$ and~$C_{\Lambda_1}$. For a subgroup $\Gamma\subseteq\Gamma^+(\lambda_0)$ of finite index,
its \emph{Tanner code} on $\Gamma\backslash\Tree$ is the space of $\Gamma$-invariant families of sections
$s_v\in H^0(C_v,L_v)$ whose values at the two special points of each edge correspond.

By \cref{lem:uniformization}, the lines of $\mathcal M^{(n)}_{\mathrm{red}}$ with the restrictions of~$L$, their
superspecial points and the identifications of the fibres of $L$ at the crossings form a geometric datum $D_{\mathrm{geo}}$
on~$\Tree$, the same for every~$k$. The stabilizers act faithfully on the lines: the stabilizer of a vertex acts on its
line through its action on $\Lambda/3\Lambda$, or, for a line of the second kind, through its action on
$\Phi\Lambda'/3\Phi\Lambda'$ conjugated by~$\sigma$, since $J$ commutes with~$\Phi$; and $2I\to\SL_2(\F_9)$ is injective (proof
of \cref{lem:amalgam}).

\begin{lemma}[Sections are Tanner codes]\label{lem:sections}
\notformalized
For every subgroup $\Gamma\subseteq\Gamma^+(\lambda_0)$ of finite index, $H^0(\Gamma\backslash\mathcal M^{(n)}_{\mathrm{red}},L)$
is isomorphic to the Tanner code of $D_{\mathrm{geo}}$ on $\Gamma\backslash\Tree$.
\end{lemma}
\begin{informalproof}
Since $\Gamma$ acts freely on $\Tree$ (\cref{lem:chain}), the curve $Y=\Gamma\backslash\mathcal M^{(n)}_{\mathrm{red}}$ has
the lines as components and ordinary nodes exactly at the superspecial points. So
$0\to\OO_Y\to\nu_*\OO_{\tilde Y}\to\bigoplus_{\text{nodes}}\bar\F_3\to0$ for the normalization~$\nu$, and a section of $L$
on $Y$ is a family of sections on the components whose values at each node correspond, that is, an element of the Tanner
code.
\end{informalproof}

\subsection{Rigidity and the twist}\label{sec:rigidity}

Let $\rho\colon2I\to\SL_2(\F_9)$ be reduction modulo~$3$ through $I\otimes_{\OO_E}\OO_{E,3}\cong M_2(\Z_9)$, that is, the
action of $2I$ on $\Lambda_0/3\Lambda_0=\F_9^2$. It is one of the two representations of \cref{lem:A5}(a), and the other
is $\rho^\sigma$, obtained by applying $\sigma$ to the entries. Applying $\sigma$ to the coefficients turns a datum into
another one, whose Tanner codes are the images of the old ones under~$\sigma$, so of the same dimension.

The data $D_{\mathrm{untw}}$ and $D_{\mathrm{tw}}$ of \cref{sec:datum} are geometric data: at every vertex take $\Pone$
with $\OO(4)$, whose sections are the binary quartic forms; take as special points the points $[x]$ of the labels~$x$; and
at an edge with labels $x$ and $y$ at its two ends, identify the fibres so that sections $f$ and $g$ correspond when
$f(x)=g(y)$, which is well defined because labels are defined up to scalars $c$ with $c^4=1$. Since evaluation of quartics
at ten points is injective (\cref{lem:selfdual}), the Tanner code of $D_{\mathrm{tw}}$ on $\Gamma\backslash\Tree$ in this
sense is $T_{\mathrm{tw}}(\Gamma\backslash\Tree)\otimes_{\F_9}\bar\F_3$. In both data the stabilizer of $\Lambda_0$ acts
through~$\rho$, since $u\in2I$ maps the edge $u'e_0$, with label $u'p_0$, to the edge $uu'e_0$, with label $\rho(u)u'p_0$.
In the coordinates transported by~$\alpha$, the stabilizer of $\Lambda_1$ acts through $\rho\circ\Ad\alpha^{-1}$ in
$D_{\mathrm{untw}}$, and through $(\rho\circ\Ad\alpha^{-1})^\sigma$ in $D_{\mathrm{tw}}$, since $s=\alpha s'\alpha^{-1}$ maps the
label $\Frob(t'k_0)$ to $\Frob(s't'k_0)=\Frob(s')\Frob(t'k_0)$. So $D_{\mathrm{untw}}$ has type $(\rho,\rho\circ\Ad\alpha^{-1})$
and $D_{\mathrm{tw}}$ has type $(\rho,(\rho\circ\Ad\alpha^{-1})^\sigma)$.

\begin{lemma}[Rigidity]\label{lem:rigidity}
\notformalized
Let $D$ and $D'$ be geometric data on~$\Tree$.
\begin{enumerate}[label=(\alph*)]
\item If $D$ and $D'$ have the same type, their Tanner codes on $\Gamma\backslash\Tree$ are isomorphic for every subgroup
$\Gamma\subseteq\Gamma^+(\lambda_0)$ of finite index.
\item Every geometric datum has the type of one of $D_{\mathrm{untw}}$, $D_{\mathrm{tw}}$, $D_{\mathrm{untw}}^\sigma$
and~$D_{\mathrm{tw}}^\sigma$, and these four types are different.
\end{enumerate}
\end{lemma}
\begin{informalproof}
(1) By \cref{lem:A5}(a), the stabilizer $A_5$ of $\Lambda_0$ acts on $C_{\Lambda_0}$ through $\rho$ or $\rho^\sigma$, up to
isomorphism, and $\alpha A_5\alpha^{-1}$ acts on $C_{\Lambda_1}$ through $\rho\circ\Ad\alpha^{-1}$ or its $\sigma$-twist. So
there are at most four types, and since $\rho\not\cong\rho^\sigma$, the four data in~(b) have four different ones. This
is~(b).

(2) By \cref{lem:A5}(b), an equivariant isomorphism between two such vertex curves is unique. An $A_5$-linearization of
$\OO(4)$ is unique because $A_5$ is perfect, and its equivariant automorphisms are scalars.

(3) The ten special points at a vertex are permuted by its stabilizer, transitively with stabilizer $S_3$
(\cref{lem:amalgam}), so by \cref{lem:A5}(c) they form the orbit of size~$10$; since $S_3$ is self-normalizing, the
equivariant bijection between the edges at a vertex and its special points is unique.

(4) Take $D$ and $D'$ of the same type, fix equivariant isomorphisms of their vertex data at $\Lambda_0$ and $\Lambda_1$,
unique up to scalars by~(2) and~(3), and transport them by~$\Gamma^+$. The edge stabilizers act on the fibres at the
special points through the same character, the inverse fourth power of the character of $\Gamma_e$ on the line of the
edge, which is trivial (\cref{sec:datum}). So the ratio of the two identifications at an edge is a $\Gamma^+$-invariant
function on edges, hence a constant $r$ by edge-transitivity (\cref{lem:amalgam}), and scaling the sections at the odd
vertices by $r$ is an isomorphism of Tanner codes on every $\Gamma\backslash\Tree$. This is~(a).
\end{informalproof}

\begin{lemma}[The twisted type]\label{lem:twist}
\notformalized
$D_{\mathrm{geo}}$ has the type of $D_{\mathrm{tw}}$ or of $D_{\mathrm{tw}}^\sigma$. Hence
$h^0(Y_k,L|_{Y_k})=\dim_{\F_9}T_{\mathrm{tw}}(X_k)$ for every $k\ge0$.
\end{lemma}
\begin{informalproof}
We use the local picture of \cref{sec:uniformization} and the identification $U_0\cong\Q_9^2$ of
\cref{lem:uniformization}.

(1) An element $g\in J$ commutes with~$\Phi$. So if $b_1,b_2$ is a basis of $U_0$ and $U_1$ is given the basis
$\Phi^{-1}b_1,\Phi^{-1}b_2$, the matrix of $g$ on $U_1$ is $\sigma^{-1}=\sigma$ applied to its matrix on~$U_0$.

(2) Hence the stabilizer of a vertex acts on a line $C_\Lambda$ of the first kind through its action on $\Lambda/3\Lambda$,
and on a line $C_{\Lambda'}$ of the second kind through its action on $M/3M$, for $M=\Phi\Lambda'$, with the entries
conjugated by~$\sigma$: $\Phi$ identifies $C_{\Lambda'}$ with the Frobenius twist of $\Pone(M/3M)$.

(3) Suppose that the lines at the even vertices are of the first kind. Then $A_5$ acts on $C_{\Lambda_0}$ through~$\rho$, and
$\alpha A_5\alpha^{-1}$ acts on $C_{\Lambda_1}$ through its action on $\alpha\Lambda_0/3\alpha\Lambda_0$ conjugated by~$\sigma$,
which in the coordinates transported by~$\alpha$ is $(\rho\circ\Ad\alpha^{-1})^\sigma$. This is the type of~$D_{\mathrm{tw}}$.
The special points match the labels of \cref{sec:datum}. Up to the same power of~$3$, the lattices at $\Lambda_0$ and
$\Lambda_1$ are $\Lambda_0$ and $M=\alpha\Lambda_0$. On $C_{\Lambda_0}$ the special point towards $\Lambda_1$ is
$\alpha\Lambda_0/3\Lambda_0=\im(\alpha\bmod3)=\langle p_0\rangle$. On $C_{\Lambda_1}$ the special point towards $\Lambda_0$ is
$\Phi^{-1}(3\Lambda_0)/3\Lambda'$, which $\Phi$ maps to $3\Lambda_0/3\alpha\Lambda_0$; in the coordinates transported by
$\alpha$ this is $\bar\alpha\Lambda_0/3\Lambda_0=\ker(\alpha\bmod3)=\langle k_0\rangle$, so the special point is
$[\Frob(k_0)]$.

(4) If the lines at the even vertices are of the second kind, the same computation gives the type
$(\rho^\sigma,\rho\circ\Ad\alpha^{-1})$, which is that of~$D_{\mathrm{tw}}^\sigma$.

(5) If the embedding of $\OO_E\otimes\Z_3$ into $W$ that defines $L_0$ reduces to $\theta\mapsto2-i$, both entries of the
type are replaced by their images under~$\sigma$. In every case the type is that of $D_{\mathrm{tw}}$ or
of~$D_{\mathrm{tw}}^\sigma$.

For the last claim, $Y_k\cong\Gamma^+(\lambda_0^{k+1})\backslash\mathcal M^{(n)}_{\mathrm{red}}$ (\cref{lem:uniformization}),
so by \cref{lem:sections,lem:rigidity} $h^0(Y_k,L|_{Y_k})$ is the dimension over $\bar\F_3$ of the Tanner code on $X_k$ of
$D_{\mathrm{tw}}$ or of $D_{\mathrm{tw}}^\sigma$. This is $\dim_{\F_9}T_{\mathrm{tw}}(X_k)$: the first code is
$T_{\mathrm{tw}}(X_k)\otimes_{\F_9}\bar\F_3$, and the second its image under~$\sigma$.
\end{informalproof}

\subsection{Conclusion}\label{sec:proofH}

\begin{theorem}[Theorem H]\label{thm:H}
\notformalized
For the chain $N_k=\psi^{-1}(\Gamma^+(\lambda_0^{k+1}))$ of $F$, with $\abs{F:N_k}=11^{3k}$, and every $k\ge0$,
\[
  \dim_{\F_3}\ker\varphi_{F/N_k}(A)\ \ge\ 22\,\abs{F:N_k}-2\ \ge\ 20\,\abs{F:N_k}.
\]
In particular hypothesis (KD) holds with $\eps=20$, and so does (KN).
\end{theorem}
\begin{informalproof}
By \cref{eq:H1}, \cref{lem:twist} and \cref{lem:RR},
\[
  \dim_{\F_3}\ker\varphi_{F/N_k}(A)=2\dim_{\F_9}T_{\mathrm{tw}}(X_k)=2h^0(Y_k,L|_{Y_k})\ge2(V_k-1)=22\,\abs{F:N_k}-2,
\]
since $V_k=11^{3k+1}=11\,\abs{F:N_k}$; and $22\abs{F:N_k}-2\ge20\abs{F:N_k}$ because $\abs{F:N_k}\ge1$.
\end{informalproof}

At $k=0$ the last inequality is an equality, $20=22-2$, while the actual kernel dimension is $42$ (\cref{rem:base}).

\subsection{Numerical checks}\label{sec:numerics}

\begin{remark}[Numerical]\label{rem:numerics}
The script \path{tools/lambda_level1.py} computes the first nested level $k=1$ exactly. It lifts the coset
representatives to $\PSL_2(\Z/121)$, obtains voltages in $\ker(\PSL_2(\Z/121)\to\PSL_2(\F_{11}))\cong\F_{11}^3$, and sums the
kernels of the $1331$ twisted $55\times55$ operators over $\F_{3^{10}}$. With $c_k$ the number of cusps
($c_0=12$, $c_1=132$):
\begin{center}
\begin{tabular}{@{}cccccccc@{}}
\toprule
$k$ & $\abs{F:N_k}$ & $V_k$ & $c_k$ & $\dim T_{\mathrm{tw}}(X_k)$ & $V_k+c_k-2$ & $V_k-1$ & $\dim T_{\mathrm{untw}}(X_k)$\\
\midrule
$0$ & $1$ & $11$ & $12$ & $21$ & $21$ & $10$ & $20$\\
$1$ & $1331$ & $14641$ & $132$ & $14771$ & $14771$ & $14640$ & $2460$\\
\bottomrule
\end{tabular}
\end{center}
So the exact formula $V_k+c_k-2$ holds at $k=0,1$. At $k=1$ the untwisted datum falls below the bound of \cref{lem:RR},
which confirms \cref{lem:twist} independently, by elimination between the types of \cref{lem:rigidity}. By
\cref{lem:tanner,lem:covering} the value $14771$ means $\dim_{\F_3}\ker\varphi_{F/N_1}(A)=29542$, a normalized kernel of
about $22.2$. Each of the $1330$ nontrivial characters of $\F_{11}^3$ contributes a twisted kernel of dimension $11$ or~$12$.
With random voltages in $\F_{11}^3$ on the same graph, which define connected covers of the same degree that are not
congruence covers (\cref{rem:congruence}), only the kernel of dimension $21$ pulled back from $X_0$ survives
(\path{tools/lambda_level1.py random s}, $s=1,2,3$). The finite group theory of \cref{lem:A5} is also checked by
exhaustive computation in \path{tools/a5_checks.py}.
\end{remark}
